% ============================================================================
% CHAPTER 11: CONCLUSION
% ============================================================================

\chapter{Conclusion}
\label{ch:conclusion}

\section{Summary of Contributions}

This work has presented a comprehensive \textbf{Contract-Based Adaptive UAV Control System} that advances the state-of-the-art in autonomous drone safety through the principled application of formal methods.

\subsection{Technical Contributions}

\begin{enumerate}[label=\textbf{C\arabic*:}]
    \item \textbf{Hierarchical Contract Framework:} We developed a complete specification of Assume-Guarantee contracts for all system components---from sensors through actuators---enabling compositional verification of system-level properties without requiring analysis of internal implementations.

    \item \textbf{Predictive Horizon Planning:} We introduced a novel use of contract composition over a planning horizon using the Pacti library. This enables the system to anticipate safety violations \textit{before} they occur, providing 1--2 seconds of advance warning compared to reactive approaches.

    \item \textbf{Principled Controller Switching:} We demonstrated mathematically justified transitions between PID and H-infinity controllers based on formal contract satisfaction, replacing ad-hoc heuristic thresholds with provable switching conditions.

    \item \textbf{Multi-Layer Safety Architecture:} We implemented a defense-in-depth approach with:
    \begin{itemize}
        \item Planning layer: Predictive contract monitoring
        \item Supervision layer: Mode-based protection
        \item Control layer: CBF-based hard constraint enforcement
    \end{itemize}

    \item \textbf{Pre-flight Verification:} We enabled mathematical proof of mission feasibility before takeoff through contract composition, allowing operators to verify that a mission can be completed safely under expected conditions.
\end{enumerate}

\subsection{Implementation Achievements}

\begin{table}[htbp]
    \centering
    \caption{Implementation summary}
    \label{tab:implementation_summary}
    \begin{tabular}{@{}ll@{}}
        \toprule
        \textbf{Component} & \textbf{Status} \\
        \midrule
        Contract-aware EKF & Implemented, tested \\
        PID Controller with contracts & Implemented, tuned \\
        H-$\infty$ Controller with contracts & Implemented, validated \\
        Flight Mode Supervisor & Implemented, 4 modes \\
        Horizon Planner (Pacti) & Implemented, 3s horizon \\
        CBF Safety Filter & Implemented, 3 constraints \\
        Controller Switcher & Implemented, bumpless transfer \\
        Adaptive Control System & Integrated, demonstrated \\
        \bottomrule
    \end{tabular}
\end{table}

\section{Key Findings}

\subsection{Contract-Based Design Benefits}

\begin{enumerate}
    \item \textbf{Modularity:} Contracts provide clean interfaces between components, enabling independent development and testing.

    \item \textbf{Compositionality:} System-level properties can be verified by composing component contracts, without analyzing implementations.

    \item \textbf{Transparency:} Contract violations are explicit and traceable, supporting debugging and certification.

    \item \textbf{Adaptability:} Contracts define operating envelopes that enable principled adaptation.
\end{enumerate}

\subsection{Experimental Validation}

The multi-waypoint mission experiments demonstrated:

\begin{itemize}
    \item \textbf{100\% waypoint completion} under varying wind conditions
    \item \textbf{Zero safety violations} with CBF enforcement
    \item \textbf{Correct controller switching} based on contract satisfaction
    \item \textbf{8\% energy savings} compared to always-robust operation
    \item \textbf{Predictive advantage} of 1.5s average warning time
\end{itemize}

\subsection{Lessons Learned}

\begin{enumerate}
    \item \textbf{EKF Predict Step:} The predict step must always precede update to prevent state drift. This was a critical bug that caused significant issues until identified.

    \item \textbf{NED Coordinate System:} The NED convention (negative Z = altitude) requires careful attention throughout the codebase to avoid sign errors in altitude control.

    \item \textbf{Wind Estimation:} Direct wind sensor measurement is more reliable than drift-based estimation. Low-pass filtering ($\alpha = 0.15$) balances responsiveness and stability.

    \item \textbf{Waypoint Detection:} Confirmation counting (5 consecutive checks) prevents false waypoint-reached triggers during oscillations.

    \item \textbf{Controller Tuning:} Conservative PID gains prevent overshoot in simulation; real hardware may require different tuning.
\end{enumerate}

\section{Limitations}

\subsection{Current Limitations}

\begin{enumerate}
    \item \textbf{Simulation Only:} All results are from simulation; no flight test validation has been performed.

    \item \textbf{Simplified Dynamics:} The simulation uses simplified quadrotor dynamics; real vehicles have complex aerodynamic effects.

    \item \textbf{No Obstacle Avoidance:} The current system does not include obstacle detection or avoidance.

    \item \textbf{Static Contracts:} Contract parameters are fixed at design time; runtime contract learning is not implemented.

    \item \textbf{Computational Cost:} Pacti contract composition adds computational overhead compared to simple threshold checks.
\end{enumerate}

\subsection{Scope Boundaries}

\begin{itemize}
    \item Single vehicle only (no multi-UAV coordination)
    \item Position-based control (no visual servoing)
    \item Known environment (no SLAM)
    \item Hover-capable vehicles (no fixed-wing)
\end{itemize}

\section{Future Work}

\subsection{Near-Term Extensions}

\begin{enumerate}
    \item \textbf{Hardware-in-the-Loop Testing:}
    Interface with PX4 SITL and Gazebo for realistic simulation.

    \item \textbf{Obstacle Avoidance:}
    Integrate CBF-based obstacle avoidance with safety contracts.

    \item \textbf{Wind Model Improvement:}
    Implement Dryden or von K\'arm\'an wind models for realistic disturbances.

    \item \textbf{Performance Optimization:}
    Optimize Pacti calls and implement caching for real-time performance.
\end{enumerate}

\subsection{Medium-Term Research}

\begin{enumerate}
    \item \textbf{Learning-Based Contract Adaptation:}
    Use online learning to adapt contract parameters based on observed system behavior.

    \item \textbf{Multi-UAV Coordination:}
    Extend contracts to handle inter-vehicle interactions and collision avoidance.

    \item \textbf{Degraded Sensor Contracts:}
    Develop contracts for graceful degradation under sensor failures.

    \item \textbf{Formal Verification Tools:}
    Integrate with model checkers (e.g., UPPAAL, PRISM) for stronger guarantees.
\end{enumerate}

\subsection{Long-Term Vision}

\begin{enumerate}
    \item \textbf{Certification Framework:}
    Develop a complete certification package for DO-178C compliance using contract-based evidence.

    \item \textbf{Commercial Integration:}
    Partner with UAV manufacturers to integrate contract-based control into commercial products.

    \item \textbf{Urban Air Mobility:}
    Adapt the framework for passenger-carrying autonomous aerial vehicles.

    \item \textbf{General Autonomous Systems:}
    Extend the contract-based approach to other autonomous systems (ground vehicles, robots, etc.).
\end{enumerate}

\section{Broader Impact}

\subsection{Scientific Contribution}

This work demonstrates the practical applicability of formal methods to real-time control systems. The combination of:
\begin{itemize}
    \item Assume-Guarantee contracts
    \item Polyhedral contract algebra (Pacti)
    \item Control Barrier Functions
    \item Robust control (H-$\infty$)
\end{itemize}
creates a framework that bridges the gap between formal methods research and practical engineering.

\subsection{Societal Impact}

As autonomous drones become more prevalent in society, the ability to provide formal safety guarantees becomes increasingly important:

\begin{itemize}
    \item \textbf{Public Safety:} Formal guarantees reduce the risk of drone-related accidents.
    \item \textbf{Regulatory Acceptance:} Contract-based evidence supports certification processes.
    \item \textbf{Transparency:} Explainable decisions build public trust in autonomous systems.
    \item \textbf{Insurance:} Formal guarantees may reduce insurance costs for commercial operators.
\end{itemize}

\section{Final Remarks}

The Contract-Based Adaptive UAV Control System represents a significant step toward provably safe autonomous flight. By combining formal methods with practical control engineering, we have demonstrated that:

\begin{quote}
\textit{Safety and adaptability are not mutually exclusive---through principled contract-based design, autonomous systems can be both formally verified and dynamically responsive to changing conditions.}
\end{quote}

As the UAV industry matures and regulatory requirements tighten, approaches like the one presented here will become essential for deploying autonomous systems in safety-critical applications. We believe this work provides a foundation for the next generation of verifiably safe autonomous aerial vehicles.

\vspace{1cm}

\begin{center}
\begin{tikzpicture}
    \node[draw, rounded corners=10pt, fill=contractpurple!10, inner sep=15pt, text width=10cm, align=center] {
        \textbf{\Large Contract-Based Adaptive UAV Control}\\[10pt]
        \textit{Formal Methods for Provably Safe Autonomous Flight}\\[15pt]
        \begin{tabular}{@{}ccc@{}}
            \includegraphics[height=1cm]{figures/placeholder_contracts.pdf} &
            \includegraphics[height=1cm]{figures/placeholder_uav.pdf} &
            \includegraphics[height=1cm]{figures/placeholder_safety.pdf} \\
            Contracts & Autonomy & Safety
        \end{tabular}
    };
\end{tikzpicture}
\end{center}

